\section{Connections to the B and C Tracks}
\label{sec:connections}

The F track is designed to complement, not replace, the B-track congressional
redistricting research. We note three primary connections.

\subsection{B.5/B.6: Computational Complexity Scaling}
\label{sec:connections:b5b6}

Papers B.5 \citep{dellaLibera2026nwayRecursive} and B.6
\citep{dellaLibera2026complexity} characterise the scaling behaviour of
recursive bisection as a function of $k$ and $n$ (number of geographic units).
F.3 applies these scaling results to state legislative redistricting, where
$k$ can exceed $n_{\mathrm{tracts}}$. The main additional challenge is that
B.5 and B.6 assume $k \ll n$; F.3 must handle the degenerate regime where
$k \approx n$ or $k > n$.

The primary new finding for state legislative applications is that runtime
scales as $O(k \cdot n \cdot \log n)$ in the normal regime but degrades
to $O(n^2 \cdot \log n)$ when $k/n$ approaches $1$, because the METIS
partitioner must maintain contiguity across many small sub-regions that share
few graph edges. The resolution upgrade (tract $\to$ block-group) restores
the normal scaling regime.

\subsection{B.13: NestSection for Bicameral Legislatures}
\label{sec:connections:b13}

Paper B.13 \citep{dellaLibera2026nestSection} introduced NestSection for the
federal context: drawing both the U.S.\ House map and a hypothetical 100-district
Senate map as nested geographic hierarchies. F.2 applies NestSection directly
to the bicameral state legislature context, where the nesting requirement is
constitutionally mandated (in many states, senate districts must be exact unions
of house districts or of county units).

The key extension is that state bicameral ratios are far more varied than the
federal case. The federal $H:S$ ratio is $435:100 = 87:20$, with
$\gcd(435,100) = 5$---a 5-district spine. State ratios include $2:1$ (most
common, e.g., WA $98:49$, MN $134:67$), $3:1$ (OH $99:33$), $4:1$ (TN $99:33$
... correction: TN $99:33$ is $3:1$), and many irregular ratios with small
$\gcd$ values. F.2 catalogues all 49 bicameral ratios and their NestSection
compatibility.

\subsection{C.1: MAUP Sensitivity at Sub-Congressional Scale}
\label{sec:connections:c1}

Paper C.1 \citep{dellaLibera2026maup} analysed the Modifiable Areal Unit
Problem (MAUP) for congressional redistricting, comparing outcomes at tract,
block-group, and block resolutions. F.3 extends this analysis to the state
legislative regime, where the MAUP is more acute: the choice of base resolution
changes not only the boundary positions but also the \emph{feasibility} of the
redistricting problem (when $k > n_{\mathrm{tracts}}$, tract-resolution
redistricting is infeasible).

The new finding in F.3 is that the MAUP effect on partisan outcomes is larger
for state legislative redistricting than for congressional redistricting.
At the congressional level, C.1 found that resolution choice changes partisan
outcomes by at most $\pm 1$ seat (for a $k=52$ state). At the state legislative
level, F.3 finds that resolution choice changes outcomes by $\pm 2$--$4$ seats
for $k=98$ (Washington House), because the smaller district size amplifies
sensitivity to local population concentrations.
